\section{Background}
\label{sec:background}

\subsection{Group Relative Policy Optimization}
\label{sec:bg-grpo}

GRPO \citep{shao2024deepseekmath} is a variant of policy gradient methods for LLM fine-tuning that eliminates the critic network used in PPO \citep{schulman2017proximal}. For each prompt $x$, the model samples a group of $G$ responses $\{o_1, \ldots, o_G\}$ from the old policy $\pi_{\theta_{\text{old}}}$, and the objective is:
\begin{equation}
\label{eq:grpo-loss}
\resizebox{\columnwidth}{!}{$
\begin{aligned}
\mathcal{L}_{\text{GRPO}}=
\frac{1}{G}\sum_{i=1}^{G}\frac{1}{|o_i|}\sum_{t=1}^{|o_i|}
\Bigl\{
  \min\Bigl[
  \rho(\theta)\,\hat{A}_{i,t},\;
  \operatorname{clip}\!\Bigl(\rho(\theta), 1\!-\!\varepsilon,\, 1\!+\!\varepsilon\Bigr)\hat{A}_{i,t}
  \Bigr]
  -\beta\,\mathrm{D}_{\mathrm{KL}}
\Bigr\},
\end{aligned}
$}
\end{equation}
where $\rho(\theta) = \pi_\theta(o_{i,t} | x, o_{i,<t}) / \pi_{\theta_{\text{old}}}(o_{i,t} | x, o_{i,<t})$ is the importance sampling ratio. Each response $o_i$ receives a scalar reward $r_i$, and the advantage is computed via group-level normalization:
\begin{equation}
\label{eq:grpo-adv}
\hat{A}_{i,t} = \frac{r_i - \mathrm{mean}(\mathbf{r})}{\mathrm{std}(\mathbf{r})}.
\end{equation}
By replacing the learned value function with sample-based normalization, GRPO avoids training a separate critic model, substantially reducing memory and compute requirements.

\subsection{Reward Models for Mathematical Reasoning}
\label{sec:bg-rm}

\paragraph{Outcome Reward Models (ORM).}
An ORM provides a binary signal based solely on the final answer: $r^{\text{out}} = \mathbf{1}[\hat{a} = a^*]$, where $\hat{a}$ is the predicted answer and $a^*$ is the ground truth. For mathematical reasoning, ORMs are typically implemented as rule-based answer checkers that compare extracted answers against reference solutions \citep{hendrycks2021measuring}. While reliable and deterministic, ORMs provide no information about the quality of the reasoning process.

\paragraph{Process Reward Models (PRM).}
PRMs evaluate the quality of intermediate reasoning steps \citep{zheng2023judgingllmasajudgemtbenchchatbot,lightman2023lets, wang2024math}. We focus on \emph{rubric-based PRMs} implemented via the LLM-as-Judge paradigm \citep{yuan2025miraclesteps,Sheng2026ReinforcingCR,Yang2026}, where a capable LLM evaluates the solution against a scoring rubric. In our setting, a rubric-based PRM assigns a score $r^{\text{proc}} \in \{0, 0.5, 1.0\}$ reflecting whether the reasoning is fully correct, scored as 1.0, largely correct with minor issues, scored as 0.5, or fatally flawed, scored as 0.0. The complete rubric prompt is provided in Appendix~\ref{sec:rubric-prompt}.

Rubric-based PRMs offer several advantages: they require no step-level annotations, provide interpretable feedback, and can leverage the growing capabilities of frontier LLMs. However, as we demonstrate in \S\ref{sec:problem}, naive integration of PRM scores into GRPO leads to training instability.
